\chapter{Word transformations}
\MARSChapterHeading{Word transformations}
\label{revision:chacha-lab}

This laboratory accompanies Chapter 22 of the main book.
Use the complete worked program at
\texttt{revision/labs/files/chacha20.asm} to inspect a known implementation.
Use a separate copy of package \texttt{026-chacha20/starter/main.asm} for your
own construction. Keeping those two roles separate lets you compare a boundary
without replacing the program you are trying to understand.

The completed entries are \texttt{chacha\_quarter},
\texttt{chacha\_block\_scalar}, \texttt{chacha\_block\_msa} and
\texttt{chacha\_xor}. Their detailed memory and return contracts are in
\texttt{assignment.md}. The caller's labels are fixtures for supplying data;
functions must still use the addresses passed in their arguments. The completed
checks can relocate objects, change gaps and repeat calls in one session.

\section{Use a prediction to locate the first dependency}

Leave \texttt{mode} at zero in the worked program and put a breakpoint at
\texttt{after\_qr\_group1}. Before running, write the expected low 32 bits of
\texttt{t0} through \texttt{t3}. Compute the addition, then XOR, then rotation
in that order. Record the operands for the XOR as well as its result. That
extra line distinguishes using an old value from rotating in the wrong direction.

At this boundary, \texttt{t0} should be \texttt{12131415} and \texttt{t3}
should be \texttt{51721330}. The other two values have not changed. If your
register view shows additional high digits, compare the word portion first and
explain the word instruction's extension separately. Do not change an
instruction solely to make all sixteen displayed digits look unsigned.

The selected words in \texttt{quad} remain unchanged until the stores at the
end of the function. Register state and caller memory therefore answer different
questions at this breakpoint. Step to \texttt{main\_done} to inspect the four
published results, at offsets 0, 20, 40 and 60. The gaps should still contain
\texttt{a5} bytes.

\Needspace{10\baselineskip}
In your implementation, complete only the quarter-round first. Run its public
check before building the block loop:

\begin{Verbatim}[fontsize=\small]
java -jar /path/to/cli.jar grade \
  packages/026-chacha20/checks/quarter.kts \
  path/to/your/main.asm
\end{Verbatim}

Paths are relative to your book checkout in this example; replace the CLI and
source paths with your own. A passing result establishes the tested cases,
including argument selection. It does not make the other three entries complete.
The starter's status three deliberately identifies unfinished work.

\section{Separate selection errors from arithmetic errors}

Set \texttt{mode} to one to run the scalar block. The caller's supplied state
has counter one. Stop after its first quarter-round call and compare only word
positions 0, 4, 8 and 12 with a separate quarter-round calculation. The other
twelve positions should still match the copied input. If more positions changed,
investigate the selected addresses and write bounds before the repetition count.

Next, inspect one complete double-round. Write down all eight selections as
word indices, then convert each index to a byte offset. This catches a common
mistake: using position 12 as byte offset 12 instead of byte offset 48.
Do not compare a one-double-round working state with the final serialized block;
the remaining nine repetitions and feed-forward have not happened yet.

At \texttt{main\_done}, the first sixteen output bytes for the supplied block
are \texttt{10 f1 e7 e4 d1 3b 59 15 50 0f dd 1f a3 20 71 c4}.
Compare the full 64-byte block with the scalar check, rather than inferring that
one matching row establishes the rest. Also compare the original 64 input bytes:
the block entry promises to preserve them.

Implement the scalar block in your source and run \texttt{checks/scalar.kts}.
When diagnosing a mismatch, choose the earliest of these boundaries that is
wrong: copy, first quarter-round, first double-round, tenth double-round,
feed-forward, serialized bytes. Each boundary removes several possible causes.

\section{Explain a shuffle before executing it}

Set \texttt{mode} to two. At \texttt{after\_msa\_load}, each vector register
contains one row of the initial state. Write a table with columns for destination
lanes zero through three and rows for \texttt{w0} through \texttt{w3}.
Initially label its cells with word positions, not numerical word contents.

In a second table, show the arrangement after \texttt{after\_msa\_select}.
For each destination lane, list the four positions it now connects. Lane zero
must describe the first diagonal. These are provenance labels: after a column
round, the numerical contents at those positions have already changed.
Expecting the literal integers 0 through 15 in the debugger would confuse the
diagram with the data it describes.

Now decode immediate 147 into four two-bit fields, starting at the low end.
Explain why it undoes the row-one selection with immediate 57. Row two's
selection is its own inverse; show this by composing its lane map with itself.
Only then step through the actual shuffles and compare the lane contents with
your predicted rearrangement of the pre-shuffle values.

Complete the MSA block and run \texttt{checks/msa.kts}. Compare its final output
with both the scalar function and the independent expected block. Agreement
between your own two routines is useful, but a shared selection mistake can make
them agree on the same wrong result.

For a diagnostic experiment in a disposable copy, omit one inverse shuffle.
Predict whether the first column half, the first diagonal arithmetic or the next
column half first differs from the specified arrangement. Restore the operation
after observing that boundary. The experiment is about locating the dependency;
it is not an alternative cipher.

\section{Treat the wrapper as a separate program}

Set \texttt{mode} to three for the byte wrapper. With count 65, predict two
block calls and exactly 65 destination bytes. Describe the second fragment's
source read and destination write before running. There is one of each, even
though the private generated block contains 64 bytes.

First test disjoint source and destination. Then set \texttt{in\_place} to one
and rerun from the same initial source and state. Compare the transformed source
with the previous destination. The function's load-before-store order makes
these cases equivalent. A second run over already transformed bytes is a
separate experiment: with the same state it recovers the original source.

Use a zero-filled source to expose the keystream directly, because XOR with zero
preserves its other operand. This is a useful way to distinguish an incorrect
block generator from an incorrect message walker. Use a nonzero source afterward
to verify that the wrapper actually reads and combines the supplied bytes.

Run \texttt{checks/xor.kts}. Include lengths zero, one, 63, 64, 65 and 4096
in your own notes. For each positive length, calculate the number of generated
blocks and the final fragment size. Explain why 64 has one full fragment and
no empty second iteration.

\Needspace{15\baselineskip}
\section{Write the error proof before the error branch}

For each request below, state its status and whether it may read the state,
read the message or write the destination. The no-access rules for zero and
excessive counts apply even when their supplied pointers are invalid.

\begin{center}
\begin{tabular}{ll}
\textbf{Count} & \textbf{Initial counter when applicable} \\
0 & no state required \\
4097 & no state required \\
64 & \texttt{ffffffff} \\
65 & \texttt{ffffffff} \\
65 & \texttt{fffffffe}
\end{tabular}
\end{center}

For the failing counter case, fill the destination with a recognizable pattern
before the call. Confirm both the returned status and the unchanged entire
caller-owned region. A status-two result after writing 64 bytes is still a
failure of this interface. Checking only the status would miss it.

\begin{figure}[!ht]
\centering
\begin{tikzpicture}[alt={Actual light-theme native Memory view after counter exhaustion. Both pictured byte rows remain filled with a5. The outline encloses the destination prefix. The crop does not show the separately verified status register.}]
\node[anchor=north west,inner sep=0,font=\small\bfseries] (title) at (0,0) {Failure preserves the destination};
\node[anchor=north west,inner sep=0] (a) at ([yshift=-2mm]title.south west) {\MARSRevisionCrop{revision/captures/chacha/chacha-counter.png}{0.99\linewidth}{0.03125}{0.31555556}{0.58203125}{0.60888889}};
\begin{scope}[shift={(a.south west)},x={($(a.south east)-(a.south west)$)},y={($(a.north west)-(a.south west)$)}]
\draw[black,line width=.7pt] (.20,.005) rectangle (.965,.6176);
\end{scope}
\end{tikzpicture}
\caption{A 65-byte request at counter ffffffff returns status two before writing. The first 32 destination bytes remain a5 in this crop; the separate boundary check verifies status and all 4160 destination bytes. The status register is not pictured.}
\label{fig:revision-chacha-native-counter}
\end{figure}


Repeat independent calls with the same state and different buffers, then with
changed state at the same address. Run \texttt{checks/reuse.kts} and
\texttt{checks/caller.kts}. Write why a hidden counter or a cached pointer would
be wrong even if one invocation produced the expected ciphertext.

\section{Finish with an implementation you can explain}

Guided construction retains the package's named regions and illustrated
allocation. Free writing can use another register allocation, frame or
decomposition while preserving the four public interfaces. Follow your
instructor's requirements for checks and hints.

For an independent attempt, close the worked source and implement the four
entries from the contract and your boundary notes. Start with the quarter-round,
then the scalar block, then MSA, then the wrapper. Use the six public checks as
feedback. A full passing run should be accompanied by a short explanation of
why each selected memory interval is legal and which values survive nested calls.

\Needspace{9\baselineskip}
After that attempt, design a continuation interface on paper. Specify where it
keeps the unused keystream bytes, how many remain, and which counter the next
block would use. Trace calls of lengths one and 64. They must together match a
single 65-byte transformation when started at the same stream position.
Do not silently change \texttt{chacha\_xor}: its restart semantics remain part
of the completed assignment. A continuation extension needs its own contract
and tests before it can be treated as a supported entry.

When comparing scalar and MSA work, use identical input, output and stopping
boundaries. Report a simulator instruction count as instructions, and host
elapsed time as time. Neither a passing vector nor a lower instruction count
establishes a hardware timing or security property. The result you can defend
here is an implementation of the stated transformation and memory contract.

\section{Companion lab: compute a SHA-256 fingerprint}
\label{revision-lab:sha256}

A ChaCha20 block generates bytes to combine with a message. SHA-256 instead
reduces a message to a fixed 32-byte \emph{digest}, also called a fingerprint.
The calculation must depend on every input byte and on where the message ends.
A useful implementation therefore needs more than a compression loop: it needs
a representation of unfinished input and an explicit finalization operation.
This companion revises original project \texttt{025-sha256}.

Use the lesson package \emph{Carry a message through SHA-256} alongside this
companion. Its short worked programs expose particular boundaries; you complete
the six entries in one continuing \texttt{independent/sha256.asm} source. The book checkout contains the same starter at
\path{revision/labs/files/sha256.asm}.
It supplies the caller, data arenas and constants; each unfinished entry stops
with \texttt{break}. Choose your own private labels, temporaries, helpers and
frames under the public contracts. Functions, file ownership and word
transformations are prerequisites.

Answer the inline predictions, then set a breakpoint on the requested instruction
and Debug that source. The selected paused frame and blue execution row identify
where the program stopped. After accepting an observation, the lesson may already
point to the next instruction in yellow. Move the breakpoint to that next target
and use Resume in the same process. Finish the process after the final stop,
before launching another worked program.

\begin{figure}[!ht]
\centering
\begin{tikzpicture}[alt={Three crops from one actual light-theme IDE photograph. The blue execution row is the store at schedule.asm line 27, confirmed by the selected paused frame. Register t0 contains hexadecimal 0000000061626380. The yellow lesson pointer has already moved to words_ready at line 34, the next requested stop. Label A identifies the current paused instruction; label B identifies the next lesson target.}]
\node[anchor=north west,inner sep=0,font=\small\bfseries] (title) at (0,0) {Distinguish the pause from the next lesson target};
\node[anchor=north west,inner sep=0] (code) at ([yshift=-2mm]title.south west) {\MARSRevisionCrop{revision/captures/sha/sha-native-pause.png}{0.96\linewidth}{0.42968750}{0.53666667}{0.17968750}{0.23888889}};
\begin{scope}[shift={(code.south west)},x={($(code.south east)-(code.south west)$)},y={($(code.north west)-(code.south west)$)}]
\draw[black,line width=.7pt] (.11,.7723) rectangle (.995,.8960);
\draw[black,line width=.7pt] (.11,.01485) rectangle (.995,.13366);
\node[draw=black,fill=white,inner sep=1pt,font=\small\bfseries] at (.035,.84) {A};
\node[draw=black,fill=white,inner sep=1pt,font=\small\bfseries] at (.035,.075) {B};
\end{scope}
\node[anchor=north west,inner sep=0] (frame) at ([yshift=-2mm]code.south west) {\MARSRevisionCrop{revision/captures/sha/sha-native-pause.png}{0.4416\linewidth}{0.03750000}{0.21555556}{0.78281250}{0.74888889}};
\node[anchor=north west,inner sep=0] (value) at ([yshift=-2mm]frame.south west) {\MARSRevisionCrop{revision/captures/sha/sha-native-pause.png}{0.6528\linewidth}{0.24765625}{0.07666667}{0.48671875}{0.89777778}};
\end{tikzpicture}
\caption{The selected frame confirms the current pause at line 27 (A), before
\texttt{sw} stores numeric word \texttt{61626380}. The lesson has accepted that
observation and points to line 34 (B). Move the breakpoint to the next requested
instruction and use Resume. The frame and register crops belong to the same
photograph; the program has not advanced to the yellow target.}
\label{fig:revision-sha-native-pause}
\end{figure}


The implementation checks follow the dependencies: schedule, compression,
initialization, update, finalization and file ownership. The update-only checks
inspect a returned context before finalization is available. Construction edits
can invalidate earlier results, so the last assessment rechecks all original
cases on the completed source. Keep a reasoning argument beside those results;
a finite test grid does not explain why another input should work.

The algorithm is specified in
\href{https://doi.org/10.6028/NIST.FIPS.180-4}{FIPS 180-4}.
\href{https://www.rfc-editor.org/rfc/rfc6234.html}{RFC 6234} provides additional
implementation and test material. Use those references for the fixed definitions,
not as a substitute for tracing your own state transitions. The local assignment
specifies our context layout, bounded calls and simulator file interface.

\subsection*{Distinguish input bytes, numeric words and printed digits}

Start with the three ASCII bytes for \texttt{abc}: hexadecimal
\texttt{61 62 63}. No terminating zero or newline is part of this message.
Its digest has 32 bytes. Printing each byte as two hexadecimal digits produces
64 characters; those characters are a representation of the digest, not another
32-byte message to hash.

Internally, SHA-256 processes 64-byte blocks. Its first sixteen schedule words
interpret successive groups of four block bytes in big-endian order.
The bytes \texttt{61 62 63 80} therefore represent numeric word
\texttt{61626380}. Loading that group with a native little-endian word load
would instead produce \texttt{80636261}. Decode the bytes deliberately, even
though the block happens to lie at an aligned address in one example.
The interface also accepts unaligned block starts.

After decoding, the schedule is an array of \emph{numeric} words. A normal word
store in this simulator represents \texttt{61626380} in memory as
\texttt{80 63 62 61}. That storage order is correct for the schedule array.
The final digest has a different interface: it must be written as big-endian
bytes. Decide what an object represents before deciding whether its byte order
is wrong. A byte view and a word view can both be correct.

Implement \texttt{sha256\_schedule(a0 block,a1 schedule)} first. It reads
exactly 64 block bytes and replaces all 64 words of a word-aligned, 256-byte
schedule. Previous workspace contents must not affect the result. Source and
output are disjoint. Run \texttt{checks/kernel.kts} while developing; failures
for entries you have not implemented are expected. Use the individual case names
to distinguish a completed schedule from incomplete compression or streaming.

\subsection*{Expand a schedule with four explicit predecessors}

The first sixteen words come directly from the block. For indices 16 through
63, compute a wrapping word sum:
\[
 W_t=W_{t-16}+\sigma_0(W_{t-15})+W_{t-7}+\sigma_1(W_{t-2})
 \pmod {2^{32}}.
\]
The two small-sigma functions combine rotations and a logical shift:
\[
\begin{aligned}
\sigma_0(x)&=\operatorname{rotr}(x,7)\mathbin{\oplus}
 \operatorname{rotr}(x,18)\mathbin{\oplus}(x\gg3),\\
\sigma_1(x)&=\operatorname{rotr}(x,17)\mathbin{\oplus}
 \operatorname{rotr}(x,19)\mathbin{\oplus}(x\gg10).
\end{aligned}
\]
Here the right shifts insert zeros. An arithmetic shift would copy the sign bit
and change the function for words whose high bit is one. All rotations and sums
operate on 32-bit words, just as in the quarter-round chapter.

If a pointer names the destination for $W_t$, then its four predecessors are
at byte offsets $-64$, $-60$, $-28$ and $-8$. Derive each offset by multiplying
the index difference by four. This gives a way to debug an incorrect recurrence
without guessing at constants.

\MARSCodeLabel{One possible allocation for the first small-sigma term}
\begin{minted}{mars-idea}
    lw $t2, -60($t0)
    rotr $t3, $t2, 7
    rotr $t4, $t2, 18
    xor $t3, $t3, $t4
    srl $t4, $t2, 3
    xor $t3, $t3, $t4
\end{minted}

This fragment assumes \texttt{t0} points to the next schedule word, not to the
start of the schedule. It retains the unmodified input in \texttt{t2} until
both rotations and the shift are complete. A different allocation is equally
valid if it preserves that dependency.

For the padded \texttt{abc} block developed below, $W_0$ is
\texttt{61626380}, $W_1$ through $W_{14}$ are zero, and $W_{15}$ is
\texttt{00000018}. The next four words are \texttt{61626380},
\texttt{000f0000}, \texttt{7da86405} and \texttt{600003c6}.
Predict $W_{16}$ using its four predecessors before running. Then inspect all
four expanded words. If $W_{16}$ is correct and $W_{17}$ is wrong, examine which
predecessors changed between those two iterations.

\begin{figure}[!ht]
\centering
\begin{tikzpicture}[alt={Actual native Memory row for schedule words sixteen through nineteen. Bytes are 80 63 62 61, 00 00 0f 00, 05 64 a8 7d and c6 03 00 60. An outline identifies the first numeric word.}]
\node[anchor=north west,inner sep=0,font=\small\bfseries] (title) at (0,0) {Expanded schedule words use native storage};
\node[anchor=north west,inner sep=0] (a) at ([yshift=-2mm]title.south west) {\MARSRevisionCrop{revision/captures/sha/sha-schedule.png}{0.99\linewidth}{0.03125}{0.33888889}{0.58203125}{0.60888889}};
\begin{scope}[shift={(a.south west)},x={($(a.south east)-(a.south west)$)},y={($(a.north west)-(a.south west)$)}]
\draw[black,line width=.7pt] (.20,.01) rectangle (.39,.467);
\end{scope}
\end{tikzpicture}
\caption{The pictured row contains numeric words W16 through W19. Its first four bytes, 80 63 62 61, store word 61626380 in native little-endian order. An outline identifies that first word.}
\label{fig:revision-sha-native-schedule}
\end{figure}


\subsection*{Keep a round's old state until its new values are ready}

\texttt{sha256\_compress(a0 state,a1 block,a2 workspace)} receives eight
aligned numeric state words, one 64-byte block and a 256-byte workspace.
It expands that block's schedule, performs 64 rounds and adds the working
results to the original state words. The source block is preserved. This is one
compression operation; it does not choose the initial state or add padding.

Name the eight working words $a$ through $h$. The large-sigma functions contain
three rotations each:
\[
\begin{aligned}
\Sigma_0(x)&=\operatorname{rotr}(x,2)\mathbin{\oplus}
 \operatorname{rotr}(x,13)\mathbin{\oplus}\operatorname{rotr}(x,22),\\
\Sigma_1(x)&=\operatorname{rotr}(x,6)\mathbin{\oplus}
 \operatorname{rotr}(x,11)\mathbin{\oplus}\operatorname{rotr}(x,25).
\end{aligned}
\]
They are different functions from the small sigmas in the schedule.
Two further functions operate independently at each bit position.
\emph{Choice}, $\operatorname{Ch}(e,f,g)$, selects the bit from $f$ when the
corresponding bit of $e$ is one, and from $g$ otherwise. \emph{Majority},
$\operatorname{Maj}(a,b,c)$, is one when at least two of those three input bits
are one. They can be expressed as
$(e\mathbin{\&}f)\mathbin{\oplus}(\mathord{\sim}e\mathbin{\&}g)$ and
$(a\mathbin{\&}b)\mathbin{\oplus}(a\mathbin{\&}c)
\mathbin{\oplus}(b\mathbin{\&}c)$, respectively, taking complements within
the word.

At round $t$, form two temporary word sums from the \emph{old} working state:
\[
\begin{aligned}
T_1&=h+\Sigma_1(e)+\operatorname{Ch}(e,f,g)+K_t+W_t,\\
T_2&=\Sigma_0(a)+\operatorname{Maj}(a,b,c).
\end{aligned}
\]
The supplied table gives the fixed $K_t$ values; do not derive them from message
bytes. The new working sequence is
$(T_1+T_2,a,b,c,d+T_1,e,f,g)$, with all sums modulo $2^{32}$.
Both new $a$ and new $e$ need $T_1$, and new $e$ needs old $d$.

\Needspace{15\baselineskip}
For example, if \texttt{s0} through \texttt{s7} hold $a$ through $h$, and
\texttt{t0/t1} hold $T_1/T_2$, the following order avoids destroying a value
before its last use:

\MARSCodeLabel{Publish one round's eight new working words}
\begin{minted}{mars-idea}
    move $s7, $s6
    move $s6, $s5
    move $s5, $s4
    addu $s4, $s3, $t0
    move $s3, $s2
    move $s2, $s1
    move $s1, $s0
    addu $s0, $t0, $t1
\end{minted}

Trace why reversing these moves fails. Moving $a$ into $b$ before preserving old
$b$ for $c$, for example, duplicates $a$ instead of shifting the old state.
Using saved registers for working words also means the function must save and
restore the caller's complete 64-bit values. The fact that the calculation uses
only words does not narrow the calling convention.

Starting with the supplied SHA-256 initial values and the padded \texttt{abc}
block, the first round gives $T_1=$\texttt{54da50e8} and
$T_2=$\texttt{08909ae5}. New $a$ is \texttt{5d6aebcd}; new $e$ is
\texttt{fa2a4622}. Use this early boundary to separate a wrong Boolean function
or constant from an error later in the loop. After round 64, add each working
word to its corresponding \emph{original} state word. Storing only the working
sequence omits feed-forward.

The next block of the same message starts from this updated chaining state.
Its rounds therefore depend on the preceding block. The independent columns
that made the ChaCha20 MSA arrangement useful are absent here. Several messages
could be processed independently, but that would require keeping a separate
context for each message.

\subsection*{Represent unfinished input in the context}

A caller should be able to supply \texttt{abc} at once or as updates containing
\texttt{a} and \texttt{bc}. The digest must be the same. Padding each update
would insert extra bytes between those chunks, so an update cannot finalize the
message merely because its current input ends.

\begin{table}[!ht]
  \centering
  \begin{tabular}{r r l}
    \toprule
    Byte offset & Bytes & Meaning                                    \\
    \midrule
    0           & 32    & Eight numeric chaining words               \\
    32          & 8     & Original message byte total                \\
    40          & 4     & Used bytes in the unfinished block         \\
    44          & 4     & Reserved                                   \\
    48          & 64    & Block buffer; only the used prefix is live \\
    112         & 256   & Schedule workspace                         \\
    \bottomrule
  \end{tabular}
  \caption{The 368-byte context. Numeric fields use the simulator's native little-endian storage. The source block and published digest have their own big-endian interpretation. Unused buffer and workspace bytes need not be zero.}
  \label{tab:revision-sha-context}
\end{table}


\texttt{sha256\_init(a0 context)} installs the eight standard initial words
and sets total and used to zero independently of old contents. It need not clear
unused storage. \texttt{sha256\_update(a0 context,a1 source,a2 count)} accepts
counts zero through 4096. Total original message bytes remain below $2^{61}$.
The context, source and other supplied objects are disjoint.

\Needspace{7\baselineskip}
On return from update, the chaining words describe all complete message blocks,
used equals total modulo 64, and the live buffer prefix contains the remaining
original bytes. These three facts are the invariant. Once used reaches 64,
compress the buffer and reset used to zero. Buffer bytes beyond the new used
prefix are stale storage, not additional input.

For a 65-byte message supplied as 63 bytes and then two, the first call leaves
used 63. The next byte completes a block; compress it. The final byte becomes
buffer byte zero, leaving total 65 and used one. Before coding, trace the source
pointer, remaining count and used value at all three points. Do not let a nested
compression call destroy the outer source walker.

An empty update reads no source bytes, even if its source pointer is deliberately
invalid. It can still inspect its valid context. Calling update on two contexts
in alternating order must preserve each one's independent state and tail.
Shared read-only constants are safe; one shared mutable buffer is not.

In \texttt{worked/fill-block.asm}, the two-byte request increases the original
byte total to 65 once, then only one byte is copied before the \texttt{block\_ready}
pause. Used is temporarily 64 and one input byte remains. The state words have
not absorbed the full block. This limited program stops at an internal boundary;
your update must perform that compression and finish the remaining fragment
before returning a valid context.

With mode zero, the supplied caller uses three original bytes. After update,
inspect total three, used three and the live prefix \texttt{61 62 63} before
finalization.

\subsection*{Prove whether padding needs one block or two}

\texttt{sha256\_final(a0 context,a1 output32)} consumes a prepared context
and publishes exactly 32 bytes. Finalization appends a single one bit, followed
by enough zero bits and the original message's 64-bit bit length. Since our
messages contain whole bytes, the marker byte is \texttt{80}. It is not a
terminator to search for in the original message; an original \texttt{80} byte
is ordinary data.

The final eight bytes of the last block hold the length, most significant byte
first. For \texttt{abc}, the marker is at buffer offset three, zeros extend
through offset 55, and offsets 56--63 contain
\texttt{00 00 00 00 00 00 00 18}. The final byte is hexadecimal 18 because
three original bytes contain 24 bits.

\begin{table}[!ht]
\centering
\begin{tabular}{r r r r}
\toprule
Original bytes & Blocks handled by update & Buffered bytes & Final blocks\\
\midrule
0 & 0 & 0 & 1\\
3 & 0 & 3 & 1\\
55 & 0 & 55 & 1\\
56 & 0 & 56 & 2\\
63 & 0 & 63 & 2\\
64 & 1 & 0 & 1\\
65 & 1 & 1 & 1\\
\bottomrule
\end{tabular}
\caption{Finalization always processes padding, including after an exact full message block. The last column counts additional compression calls during finalization, not the total calls for the message.}
\label{tab:revision-sha-padding}
\end{table}


\Needspace{9\baselineskip}
With 55 buffered bytes, the marker fits at offset 55 and the length starts at
56 in the same block. With 56 buffered bytes, the marker occupies offset 56,
where the length would need to start. Zero the rest of that block and compress
it first. Then create a second block with zeros through offset 55 and the
length in the last eight bytes. A full original 64-byte block has already been
compressed by update, but finalization still needs a padding block.

\begin{figure}[!ht]
\centering
\begin{tikzpicture}[alt={Actual native Memory view of the last thirty-two bytes of the second padding block for a fifty-six-byte message. All pictured bytes are zero except the final two, 01 c0. An outline identifies the eight-byte big-endian original bit length.}]
\node[anchor=north west,inner sep=0,font=\small\bfseries] (title) at (0,0) {The final block retains the original bit length};
\node[anchor=north west,inner sep=0] (a) at ([yshift=-2mm]title.south west) {\MARSRevisionCrop{revision/captures/sha/sha-padding.png}{0.99\linewidth}{0.03125}{0.31555556}{0.58203125}{0.60888889}};
\begin{scope}[shift={(a.south west)},x={($(a.south east)-(a.south west)$)},y={($(a.north west)-(a.south west)$)}]
\draw[black,line width=.7pt] (.58,.005) rectangle (.965,.303);
\end{scope}
\end{tikzpicture}
\caption{For 56 original bytes, the last two rows of the second padding block are shown before its compression. The outlined eight-byte field ends 01 c0, encoding 448 bits. Earlier message bytes and the marker are outside this crop.}
\label{fig:revision-sha-native-padding}
\end{figure}


\Needspace{8\baselineskip}
Compute the bit length from the original byte total using a doubleword left
shift by three. Serialize all eight bytes with full-width shifts. A word shift
cannot extract the upper half correctly. The public
\texttt{worked/padding-high-length.asm} experiment supplies a synthetic original
byte total of $2^{32}+63$ and a 63-byte unfinished tail. It constructs two final
blocks whose length field is \texttt{00 00 00 08 00 00 01 f8}. These values let
you inspect the high length bits without reading a file of that size. The
experiment constructs padding bytes only; the actual compression belongs to
your finalization function.

After the final compression, serialize each of the eight numeric state words
most significant byte first. For \texttt{abc}, the first four digest bytes are
\texttt{ba 78 16 bf}. Storing numeric word \texttt{ba7816bf} with a native
word store would reverse those four output bytes. Check the complete digest,
not only this prefix, against an independent digest implementation.

The \texttt{worked/serialize.asm} experiment is given the final numeric state
for \texttt{abc}. It demonstrates this output conversion without calculating
that state. Your compression and finalization must obtain the correct words
before applying the same byte-order rule.

\begin{figure}[!ht]
  \centering
  \begin{tikzpicture}[alt={Actual native Memory view of all thirty-two SHA-256 digest bytes for abc. The first row begins ba 78 16 bf and the second begins b0 03 61 a3. An outline identifies the first four serialized bytes.}]
    \node[anchor=north west,inner sep=0,font=\small\bfseries] (title) at (0,0) {The digest is a serialized byte sequence};
    \node[anchor=north west,inner sep=0] (a) at ([yshift=-2mm]title.south west) {\MARSRevisionCrop{revision/captures/sha/sha-digest.png}{0.99\linewidth}{0.03125}{0.31555556}{0.58203125}{0.60888889}};
    \begin{scope}[shift={(a.south west)},x={($(a.south east)-(a.south west)$)},y={($(a.north west)-(a.south west)$)}]
      \draw[black,line width=.7pt] (.20,.3088) rectangle (.39,.6176);
    \end{scope}
  \end{tikzpicture}
  \caption{All 32 digest bytes for the three-byte message abc are shown. The outlined prefix is ba 78 16 bf. Interpret the Bytes column in order; a native little-endian word view groups these digest bytes differently.}
  \label{fig:revision-sha-native-digest}
\end{figure}


\Needspace{6\baselineskip}
Finalization consumes the context. Its contents afterward are unspecified,
and calling final twice is outside the interface. Initialize it again before
hashing another message. Do not add padding bytes to the saved original-byte
count while calculating their own length field.

\subsection*{Finalize a file only after successful EOF}

The file function receives a terminated path, scratch pointer, capacity from one
through 256, and a 32-byte digest destination. It owns a fresh private context.
It opens the path read-only, initializes that context and repeatedly reads a
positive request no larger than capacity.

A positive result tells how many bytes are now available in scratch. Supply
exactly that count to update. A short positive result does not mean EOF; the next
read may produce more data. A negative result is a read error. Zero is successful
EOF, at which point finalization can publish the digest.

Suppose capacity is seven but a controlled reader returns three bytes, then two,
then minus one. The context has consumed five successful bytes, but there is no
complete file digest to publish. Return status two, close the acquired descriptor
and preserve all 32 destination bytes. Hashing the requested seven bytes would
incorporate stale scratch contents; treating minus one as EOF would publish a
digest of a prefix.

The public \texttt{worked/read-prefixes.asm} program separates file ownership
from hash arithmetic. It counts successful input bytes and publishes one count
only at EOF. Its successful lesson fixture returns two bytes, then one, then
EOF; the error fixture returns two bytes, then minus one, preserving the original
output word. The SHA wrapper adds a fresh context, updates and finalization
under the same success/failure ownership rule, applied to all 32 digest bytes.

\begin{figure}[!ht]
  \centering
  \begin{tikzpicture}[alt={Actual native Memory view after a deterministic file read error. Both rows of the thirty-two-byte digest destination remain a5. An outline encloses the unchanged destination. Status and descriptor closure are checked separately, not visible in the crop.}]
    \node[anchor=north west,inner sep=0,font=\small\bfseries] (title) at (0,0) {A failed read leaves the digest unpublished};
    \node[anchor=north west,inner sep=0] (a) at ([yshift=-2mm]title.south west) {\MARSRevisionCrop{revision/captures/sha/sha-read-error.png}{0.99\linewidth}{0.03125}{0.31555556}{0.58203125}{0.60888889}};
    \begin{scope}[shift={(a.south west)},x={($(a.south east)-(a.south west)$)},y={($(a.north west)-(a.south west)$)}]
      \draw[black,line width=.7pt] (.20,.005) rectangle (.965,.6176);
    \end{scope}
  \end{tikzpicture}
  \caption{After successful reads of three and two bytes followed by a read error, the full 32-byte destination remains a5. The outline encloses that unchanged destination. Separate checks verify status two and closing descriptor 19; those values are not pictured.}
  \label{fig:revision-sha-native-read-error}
\end{figure}


An open failure returns status one and acquires no descriptor. On success and
read failure, close the descriptor actually returned by open, which need not be
three. Keep that descriptor separate from argument registers reused by update
and finalization. Also keep the result status across the close syscall so that
its return value does not replace the function's intended status.

\Needspace{12\baselineskip}
With caller defaults retained (\texttt{mode} two, \texttt{length} three and
\texttt{buffer\_capacity} seven), prepare an ordinary host-file observation from
your completed source. From the book checkout, choose a fresh output path:

\Needspace{7\baselineskip}
\begin{Verbatim}[fontsize=\small]
python3 packages/025-sha256/prepare.py s45-1 \
  --source path/to/your/main.asm \
  --output /private/tmp/sha-file-observation.asm \
  --fixture-output /private/tmp/sha-input.bin
\end{Verbatim}

The tool refuses existing outputs. The resulting file contains deterministic
binary bytes and the source names that exact host path. An ordinary file run
checks real file integration; it does not force a short read or injected error.
The public file checks supply those separate controls.

\subsection*{Build evidence at the boundary that can fail}

Complete the independent source in place and retain the worked programs for
separate observations. For optional standalone data experiments, the preparation
tool can apply \texttt{s03-1} for the padded schedule input, \texttt{s11-1} for a
compression input with the standard initial state, \texttt{s28-2} for crossing
a buffered block boundary, and \texttt{s37-3} for two-block finalization.
These recipes install data and controls while retaining your instruction text.
Do not reuse a context after finalization just because its old buffer still
looks plausible.

Implement and test schedule, compression, initialization, update, finalization
and file ownership in that dependency order. The public completed suites are
\texttt{kernel.kts}, \texttt{file.kts}, \texttt{program.kts} and
\texttt{reuse.kts} under the package's \texttt{checks} directory. Run each with
\texttt{java -jar /path/to/cli.jar grade CHECK SOURCE}.
The caller stops at \texttt{main\_done} before storing its returned status and
at \texttt{results\_ready} after the store. These are different observation
boundaries; the later exit syscall changes \texttt{v0}.

Write a test plan covering 0, 3, 55, 56, 63, 64 and 65 bytes, then binary input
containing zeros and bytes above 127. For each length, compare one update with
several splits, including empty updates. Interleave two contexts and initialize
a used context again. Move the selected slices within their supplied arenas.
A correct digest at a fixed address does not establish correct pointer use or
memory ownership.

For error diagnosis, check the first wrong schedule word, first wrong round,
first wrong context invariant or first wrong serialized byte. A complete
fingerprint is valuable as an end-to-end check, but it cannot identify which
of those boundaries first diverged.

Finally, an unkeyed digest does not identify a sender. A party that can replace
both a message and its advertised digest can recompute them together. Keyed
message authentication requires a defined construction, such as \href{https://www.rfc-editor.org/rfc/rfc2104.html}{HMAC}; appending
a secret to this exercise is not its specification. The work here is a precise
bounded SHA-256 implementation with reproducible observations. Explain its
state and ownership promises alongside the passing results.


\section{Companion lab: extend arithmetic beyond one register}
\label{revision-lab:rsa}

Word operations have a fixed width. When a result needs more bits, an assembly
program must keep those bits explicitly. This companion develops unsigned
multiword arithmetic through a small RSA calculation. It revises original project
\texttt{027-rsa-arithmetic}; the numerical interfaces also work for composite
moduli unrelated to an RSA key.

Use the native lesson \emph{Extend arithmetic beyond one register}. Implement
the seven entries in its \texttt{independent/rsa-arithmetic.asm}. The same bare
starter is retained at \texttt{revision/labs/files/rsa-arithmetic.asm} for
standalone work. The skeleton supplies caller
objects and an aligned stack area, with a deliberate \texttt{break} at each
unfinished function. The original package retains its cumulative guided source
and public checks. Your completed functions can choose different helpers,
registers, frames and exponent scans. Start with the scalar cases before extending
the representation to four limbs.

The native lesson begins with short worked sources in its \texttt{worked}
directory. Each exposes one operation before asking you to implement a complete
interface. Answer the inline prediction, inspect the requested Debug boundaries,
then return to the continuing independent source. Finish the Debug process before
another launch or source edit. The worked files remain unchanged throughout the
required observations; there is no need to change your implementation's mode
word to follow them.

\subsection*{Use the gcd to explain a small exponent pair}

For positive modulus $N$, the \emph{residue} of $x$ is the remainder in the
range zero through $N-1$. Two integers are congruent modulo $N$ when their
difference is divisible by $N$. Reducing intermediate values to residues keeps
the numbers smaller without changing the final modular result.

Two integers are coprime when their greatest common divisor is one.
For an RSA arithmetic example, choose primes 61 and 53. Their product is
$N=3233$. The count of residues coprime to this product is
$(61-1)(53-1)=3120$. Choose public exponent 17. Its greatest common divisor
with 3120 is one, which makes an inverse modulo 3120 possible.

\begin{table}[!ht]
  \centering
  \begin{tabular}{r r r r}
    \toprule
    $A$  & $B$ & Quotient & Remainder \\
    \midrule
    3120 & 17  & 183      & 9         \\
    17   & 9   & 1        & 8         \\
    9    & 8   & 1        & 1         \\
    8    & 1   & 8        & 0         \\
    \bottomrule
  \end{tabular}
  \caption{Each Euclidean iteration replaces $(A,B)$ with $(B,A\bmod B)$. The last nonzero remainder is one. The terminating zero is not the gcd.}
  \label{tab:revision-rsa-euclid}
\end{table}


\texttt{rsa\_gcd64(a0 A,a1 B)} accepts full unsigned 64-bit values.
At each iteration, compute the unsigned remainder and replace both operands.
The loop stops when $B$ is zero, returning the current $A$. Thus
$\gcd(A,0)=A$; our local convention also defines $\gcd(0,0)=0$.
Check for a zero divisor before computing the remainder.

Begin the native observation with \texttt{worked/remainders.asm}. Its first
pair is 13,11, whose unsigned remainder is two. Before replacing either operand,
identify which old value each move still needs. The next pair must be 11,2;
its remainder is one. Figure~\ref{fig:revision-rsa-native-pause} shows that first
pause. The blue execution row and selected frame describe the current machine
state. The yellow lesson pointer identifies the next requested stop after the
lesson has accepted the current observation.

\begin{figure}[!ht]
\centering
\begin{tikzpicture}[alt={Three crops from one actual light-theme IDE photograph of the public remainders experiment. The blue execution row is the move at line 9, confirmed by the selected paused frame. Registers t0 and t1 still contain decimal 13 and 11, while t2 contains their unsigned remainder two. The yellow lesson pointer has moved to line 13, the next requested stop after updating the pair and computing another remainder. Label A identifies the current pause and label B identifies the next lesson target.}]
\node[anchor=north west,inner sep=0,font=\small\bfseries] (title) at (0,0) {Read the remainder before updating the pair};
\node[anchor=north west,inner sep=0] (code) at ([yshift=-2mm]title.south west) {\MARSRevisionCrop{revision/captures/rsa/rsa-native-pause.png}{0.96\linewidth}{0.42968750}{0.53777778}{0.17968750}{0.18888889}};
\begin{scope}[shift={(code.south west)},x={($(code.south east)-(code.south west)$)},y={($(code.north west)-(code.south west)$)}]
\draw[black,line width=.7pt] (.11,.5528) rectangle (.995,.6423);
\draw[black,line width=.7pt] (.11,.1951) rectangle (.995,.2846);
\node[draw=black,fill=white,inner sep=1pt,font=\small\bfseries] at (.075,.60) {A};
\node[draw=black,fill=white,inner sep=1pt,font=\small\bfseries] at (.075,.24) {B};
\end{scope}
\node[anchor=north west,inner sep=0] (frame) at ([yshift=-2mm]code.south west) {\MARSRevisionCrop{revision/captures/rsa/rsa-native-pause.png}{0.4416\linewidth}{0.03750000}{0.22222222}{0.78281250}{0.75111111}};
\node[anchor=north west,inner sep=0] (value) at ([yshift=-2mm]frame.south west) {\MARSRevisionCrop{revision/captures/rsa/rsa-native-pause.png}{0.768\linewidth}{0.24765625}{0.03333333}{0.43984375}{0.89777778}};
\end{tikzpicture}
\caption{At the current pause (A), \texttt{t0}=13 and \texttt{t1}=11 have
produced \texttt{t2}=2. The register view is hexadecimal: \texttt{0d} and
\texttt{0b} mean decimal 13 and 11. The move at line 9 has not executed. The lesson accepts
that observation and points to the next requested boundary at line 13 (B).
Place the next breakpoint there and use Resume: the pair becomes 11,2, and the
next unsigned remainder is one. The frame and register crops come from the same
photograph. The program is still paused at A.}
\label{fig:revision-rsa-native-pause}
\end{figure}


The final part of this worked source divides the all-one doubleword by 17.
Unsigned interpretation gives $2^{64}-1$, with remainder zero; signed
interpretation gives $-1$, with signed remainder $-1$. The register bits alone
do not choose the operation. Use the interpretation promised by the interface,
then complete the gcd assessment before moving on to modular powers.

Work backward through the 3120,17 calculation:
\[
 1=9-8=2\cdot9-17=2\cdot3120-367\cdot17.
\]
Consequently $-367$ is an inverse of 17 modulo 3120. Its nonnegative residue is
2753, and $17\cdot2753=15\cdot3120+1$. This derives the exponent used below;
it is not a second constant to memorize. You need not implement key generation
or an extended-gcd interface for this assignment.

The gcd entry returns its numerical result in \texttt{v0} and success status
zero in \texttt{v1}. Other entries return status in \texttt{v0} and write their
numerical result to memory. The supplied caller pauses at \texttt{gcd64\_return}
before moving the gcd status into \texttt{v0}. Inspect the gcd value at that
first boundary; \texttt{main\_done} is too late to expect it in \texttt{v0}.
A function returning a high-bit-set gcd must preserve all 64 result bits.

\Needspace{11\baselineskip}
\subsection*{Reduce a power without constructing the full integer}

Take message integer 65. The small forward operation is
$65^{17}\bmod3233=2790$, and the reverse calculation is
$2790^{2753}\bmod3233=65$. Euler's theorem gives
$65^{3120}\equiv1\pmod{3233}$ because 65 is coprime to 3233. Since
$17\cdot2753=1+15\cdot3120$, the combined operation is
$65(65^{3120})^{15}$, whose residue is 65.
The program will compute the remainders without ever storing $65^{17}$ itself.

\texttt{rsa\_pow32(a0 base,a1 exponent,a2 modulus,a3 output)} receives
zero-extended unsigned word operands. Its successful output is one little-endian
word. Start with result $r=1\bmod N$, current base $x=base\bmod N$ and remaining
exponent $y=exponent$. At every loop boundary,
\[
 r\,x^y\equiv \mathit{base}^{\mathit{exponent}}\pmod N.
\]
If $y$ is odd, multiply $r$ by $x$ and reduce. Shift $y$ right logically by one,
and square and reduce $x$ if another exponent bit remains. Each iteration
removes one bit, rather than performing one multiplication for each unit of the
original exponent.

\begin{table}[!ht]
  \centering
  \begin{tabular}{r r r r}
    \toprule
    Exponent before & Result after & Base after & Exponent after \\
    \midrule
    17              & 65           & 992        & 8              \\
    8               & 65           & 1232       & 4              \\
    4               & 65           & 1547       & 2              \\
    2               & 65           & 789        & 1              \\
    1               & 2790         & 789        & 0              \\
    \bottomrule
  \end{tabular}
  \caption{Least-significant-bit-first exponentiation of 65 to power 17 modulo 3233. Initially the result is one and base is 65. This version stops before squaring an unused base after the last exponent bit.}
  \label{tab:revision-rsa-power}
\end{table}


Predict the five rows before running. Explain why the result stays 65 while
remaining exponents eight, four and two are processed. Those iterations still
change the current base; skipping their squaring would make the final odd bit
multiply by the wrong power.

\begin{figure}[!ht]
\centering
\begin{tikzpicture}[alt={Actual native Memory row after scalar modular exponentiation. The first four bytes are e6 0a 00 00, encoding 2790, followed by twelve unchanged a5 bytes. An outline identifies the four-byte output.}]
\node[anchor=north west,inner sep=0,font=\small\bfseries] (title) at (0,0) {One word holds the reduced power};
\node[anchor=north west,inner sep=0] (a) at ([yshift=-2mm]title.south west) {\MARSRevisionCrop{revision/captures/rsa/rsa-word.png}{0.99\linewidth}{0.03125}{0.33888889}{0.58203125}{0.60888889}};
\begin{scope}[shift={(a.south west)},x={($(a.south east)-(a.south west)$)},y={($(a.north west)-(a.south west)$)}]
\draw[black,line width=.7pt] (.20,.01) rectangle (.39,.467);
\end{scope}
\end{tikzpicture}
\caption{The outlined bytes e6 0a 00 00 encode decimal 2790, the residue of 65 to power 17 modulo 3233. The following twelve pictured bytes remain a5 and are outside this four-byte result.}
\label{fig:revision-rsa-native-word}
\end{figure}


\Needspace{11\baselineskip}
Two word residues can have a product requiring 64 bits. Use a doubleword
multiply and an unsigned doubleword remainder for this path:

\MARSCodeLabel{Keep the entire product until after reduction}
\begin{minted}{mars-idea}
    dmulu $t1, $t1, $t0
    dmodu $t1, $t1, $a2
\end{minted}

Here \texttt{t1} is the accumulated residue, \texttt{t0} the current base and
\texttt{a2} the zero-extended modulus. A word multiply discards needed high bits.
A signed doubleword remainder also fails when the complete product's top bit is
one. For base \texttt{fffffffe}, exponent two and modulus \texttt{ffffffff},
the product is \texttt{fffffffc00000004} and the correct residue is one.
Interpreting those product bits as a negative signed integer gives a different
remainder. Use \texttt{lwu} when the caller loads unsigned word operands.

Exponent zero returns $1\bmod N$, including for a zero base. For modulus one,
that answer is zero. Modulus zero is an error: return status one without
accessing the output. Test these cases separately from a large positive exponent.

In \texttt{worked/word-product.asm}, compare \texttt{lw} and \texttt{lwu}
on the same all-one limb. Then inspect the complete product
\texttt{fffffffe00000001}. Unsigned reduction modulo \texttt{ffffffff}
gives zero, while its low word is one and signed reduction gives $-1$. The
experiment separates three values that can look plausible when printed in
hexadecimal. Finish the scalar-power checks with the complete unsigned dividend.

\subsection*{Represent a wider number with unsigned limbs}

Let $\beta=2^{32}$. Four \emph{limbs}, each an unsigned word, represent
\[
 A=A_0+A_1\beta+A_2\beta^2+A_3\beta^3.
\]
The least significant limb is at the lowest address, and the bytes inside each
limb are little-endian. This is the local arithmetic representation; it differs
from the big-endian byte strings used by the RSA conversion primitives in
\href{https://www.rfc-editor.org/rfc/rfc8017.html}{RFC 8017}.
Do not treat a dump of these limbs as an encoded RSA message.

Unsigned loads matter even when every stored limb is a perfectly valid word.
Loading \texttt{ffffffff} with \texttt{lwu} produces the value $\beta-1$ in a
64-bit register. Sign extension would instead change the upper bits used by
later doubleword multiplication. Pointer arithmetic remains full width regardless
of the limb width.

\Needspace{7\baselineskip}
\texttt{rsa\_mul128(a0 A,a1 B,a2 output)} writes the full eight-limb product.
One schoolbook implementation clears all eight output words, then processes one
input limb against each limb of the other input. For an inner product, add
$A_iB_j$, the existing output limb and the incoming carry. Store the low word
and retain the high word as the next carry.

This ordering has a useful bound. Each operand, existing limb and carry is at
most $\beta-1$, so the inner total is at most
\[
 (\beta-1)^2+(\beta-1)+(\beta-1)=\beta^2-1.
\]
It fits exactly within an unsigned doubleword. Summing several unreduced limb
products at once would require a different bound and possibly more temporary
bits. The reason this particular loop fits is its carry discipline.

For a hand-checkable example, $(2\beta-1)^2$ has low-to-high limbs
$1,\beta-4,3,0,0,0,0,0$. For the largest input $M=\beta^4-1$, its square has limbs
$1,0,0,0,\beta-2,\beta-1,\beta-1,\beta-1$. Check all eight words and the bytes immediately
after the selected 32-byte result. Returning only the low four words implements
a different operation.

\begin{figure}[!ht]
  \centering
  \begin{tikzpicture}[alt={Actual native Memory view of the complete thirty-two-byte product of two maximum unsigned 128-bit values. First-row words are one and three zeros. Second-row words are fffffffe followed by three ffffffff words. An outline encloses all eight limbs.}]
    \node[anchor=north west,inner sep=0,font=\small\bfseries] (title) at (0,0) {Keep all eight product limbs};
    \node[anchor=north west,inner sep=0] (a) at ([yshift=-2mm]title.south west) {\MARSRevisionCrop{revision/captures/rsa/rsa-product.png}{0.99\linewidth}{0.03125}{0.31555556}{0.58203125}{0.60888889}};
    \begin{scope}[shift={(a.south west)},x={($(a.south east)-(a.south west)$)},y={($(a.north west)-(a.south west)$)}]
      \draw[black,line width=.7pt] (.20,.005) rectangle (.965,.6176);
    \end{scope}
  \end{tikzpicture}
  \caption{The full product of two maximum unsigned 128-bit inputs occupies the outlined 32 bytes. Low-to-high limbs are one, zero, zero, zero, fffffffe, ffffffff, ffffffff and ffffffff.}
  \label{fig:revision-rsa-native-product}
\end{figure}


The native \texttt{worked/limb-column.asm} adds two maximum limbs and a
supplied carry of one. Its 33-bit sum is \texttt{1ffffffff}: low limb
\texttt{ffffffff}, outgoing carry one. This isolates one column. In a full
multiplication, several limb products can contribute to the same position;
account for those contributions before declaring that output limb finished.
The independent full-product checks also require zero high limbs and squaring
through one shared read-only input object.

\subsection*{Keep the carry that lies beyond the output array}

\texttt{rsa\_addmod128(a0 A,a1 B,a2 N,a3 output)} requires reduced inputs
$A<N$ and $B<N$. Their sum is below $2N$, so at most one subtraction of $N$
is needed. However, the unreduced sum can require 129 bits. Four stored result
limbs alone cannot say whether that subtraction is necessary.

Add from low limbs toward high limbs, retaining a carry into each next position.
After the fourth limb, retain the final carry separately. If it is one, subtract
$N$ even when the stored low 128 bits compare below $N$. Otherwise compare the
four-limb sum with $N$, starting at the most significant limb, and subtract when
the sum is at least $N$. Equality must produce zero.

Consider $N=2^{128}-159$ and $A=B=N-1$. The mathematical sum is
$2^{129}-320$. Its low 128 bits represent $2^{128}-320$, which is below $N$,
but its carry is one. The correct reduced result is $N-2=2^{128}-161$.
A comparison that ignores the carry would return the wrong provisional sum.
When subtracting $N$ from the low limbs, the borrow consumes the retained high
carry; the final reduced value fits in four limbs.

\begin{figure}[!ht]
  \centering
  \begin{tikzpicture}[alt={Actual native Memory row during modular addition, before conditional subtraction. The low limbs begin c0 fe ff ff and otherwise contain ff. An outline encloses these provisional sixteen bytes. The separately verified 129th carry bit is not pictured.}]
    \node[anchor=north west,inner sep=0,font=\small\bfseries] (title) at (0,0) {The provisional sum still needs subtraction};
    \node[anchor=north west,inner sep=0] (a) at ([yshift=-2mm]title.south west) {\MARSRevisionCrop{revision/captures/rsa/rsa-carry.png}{0.99\linewidth}{0.03125}{0.33888889}{0.58203125}{0.60888889}};
    \begin{scope}[shift={(a.south west)},x={($(a.south east)-(a.south west)$)},y={($(a.north west)-(a.south west)$)}]
      \draw[black,line width=.7pt] (.20,.01) rectangle (.965,.467);
    \end{scope}
  \end{tikzpicture}
  \caption{The pictured low 128 bits are the provisional sum, beginning c0 fe ff ff. Its retained high carry is one, verified separately and not pictured. Despite these low limbs being below N, subtraction must produce N minus two before return.}
  \label{fig:revision-rsa-native-carry}
\end{figure}


\Needspace{7\baselineskip}
Subtraction also proceeds from low to high, with an incoming borrow of zero
or one. For zero-extended word operands, a negative limb difference occupies
the high unsigned doubleword range; extract the borrow before retaining only
the low word. Test a borrow that propagates across several zero limbs.

This entry permits output to exactly alias either or both operands. A limbwise
implementation can load both current operand limbs before replacing the current
output limb, provided no later step needs those original limbs. The modulus
stays disjoint from output. An arbitrary partial overlap has no such guarantee;
a store could replace a higher source limb before it is read.

Use \texttt{worked/carry-reduction.asm} to inspect the decision with an
eight-bit model: $N=251$ and $A=B=250$. The full sum is 500, its low eight bits
are 244, and the outgoing carry is one. Comparing only 244 with 251 would skip
a necessary subtraction. The result is $500-251=249$. Transfer that reasoning
to the four-limb representation; the small experiment can display a complete
intermediate in one register, while your general implementation must retain its
extra carry explicitly.

\subsection*{Reduce arbitrary inputs before using residue-only helpers}

\texttt{rsa\_reduce128(a0 A,a1 N,a2 output)} accepts any unsigned 128-bit
input, not merely a value less than twice $N$. Repeatedly subtracting a small
modulus could require an impractical number of iterations. Instead, visit the
input bits from most significant to least significant.

Maintain the residue $r$ of the prefix already visited. For next bit $b$, the
new prefix has residue $(2r+b)\bmod N$. Since $r<N$, this step has a bounded
amount of reduction work. For binary input \texttt{1101}, decimal 13, modulo
five, the successive residues are one, three, one and three.
The order of the bits matters: low-to-high scanning does not follow this prefix
recurrence.

One implementation uses the reduced addition helper to compute $r+r$ and then,
when the incoming bit is one, to add one. Handle modulus one first: a literal
one is not a reduced operand for that modulus. For $N>1$, both zero and one
are valid residues. Initialize the running residue independently of any old
output or private workspace contents.

\Needspace{9\baselineskip}
\texttt{rsa\_mulmod128(a0 A,a1 B,a2 N,a3 output)} accepts arbitrary inputs
too. Reduce its current multiplicand before passing it to addmod. Keep accumulator
$r=0$, reduced current value $x=A\bmod N$ and remaining multiplier $y=B$.
When the low bit of $y$ is one, replace $r$ by $(r+x)\bmod N$. Shift $y$
right and double $x$ modulo $N$ for the next bit. The invariant is
$r+xy\equiv AB\pmod N$.

\begin{table}[!ht]
\centering
\begin{tabular}{r r r}
\toprule
Accumulator after & Current residue after & Multiplier after\\
\midrule
13 & 9 & 5\\
5 & 1 & 2\\
5 & 2 & 1\\
7 & 2 & 0\\
\bottomrule
\end{tabular}
\caption{Modular multiplication of 13 by 11 modulo 17. Initially the accumulator is zero, current residue 13 and multiplier 11. The final answer is seven; no unused last doubling is performed.}
\label{tab:revision-rsa-modproduct}
\end{table}


This computes a modular product without first constructing a 256-bit temporary
product. Both approaches can be correct, but a full-product approach also needs
a reduction algorithm that accepts all 256 product bits. Passing only the low
128 bits to \texttt{rsa\_reduce128} loses information.

Trace \texttt{worked/prefix-reduction.asm} through all four bits of
\texttt{1101} modulo five. Its returned prefix residues are 1,3,1,3. At the
third bit, the provisional six is reduced to one; the final bit then gives
three. The residue invariant explains why each append needs at most one
subtraction, while the full implementation must also retain the extra carry.
Resume to reach each new visit to the same source instruction.

In \texttt{worked/modular-product.asm}, multiplier \texttt{11} selects two
contributions from base 250 modulo 251. The first residue is 250; the doubled
base residue is 249; their sum reduces to 248. Before using the modular adder
inside your multiplier, establish its reduced-input precondition. The multiplier
itself accepts arbitrary inputs. The multiplier's caller has not promised reduced operands.

\subsection*{Shift the whole exponent and preserve the outer calculation}

\texttt{rsa\_pow128(a0 base,a1 exponent,a2 N,a3 output)} uses the same power
invariant as the scalar routine, with modular multiplication replacing the
single-register product. All 128 exponent bits participate. Scanning only its
first word or first doubleword silently changes high exponents to smaller ones.

A whole-number right shift must bring bits across storage boundaries. If two
registers hold the low and high 64-bit halves, one possible sequence is:

\MARSCodeLabel{Move the high half's low bit into the low half}
\begin{minted}{mars-idea}
    dsrl $t0, $t0, 1
    dsll $t2, $t1, 63
    or $t0, $t0, $t2
    dsrl $t1, $t1, 1
\end{minted}

For exponent $2^{64}+1$, the initial halves are one and one. After shifting,
the low half must be \texttt{8000000000000000} and the high half zero.
Shifting each half independently would instead produce a zero exponent.
For base three modulo 19, the complete exponent $2^{64}+1$ gives 13, and
$2^{127}$ gives nine. A truncated exponent can accidentally pass a test whose
true result is one, so use a case that distinguishes the two calculations.

Keep the outer result, current base, remaining exponent and modulus address
across every nested helper call. Arguments and temporary registers are volatile.
A saved return address preserves where execution resumes, not the caller's
unfinished arithmetic state. Save and restore full preserved registers, including
their upper halves, and keep local temporary arrays separate from saved addresses.

The helpers allow exact sharing among read-only inputs, so squaring may pass the
same input address twice. Modular-product and power outputs remain disjoint from
their inputs. Only addmod has the additional output-alias permission. Establish
those relationships before reusing a temporary array for another role.

The native \texttt{worked/exponent-halves.asm} performs one logical shift
of a supplied 128-bit exponent. Its old high half is
\texttt{8000000000000001} and its low half is one. After shifting, the new low
half is \texttt{8000000000000000}; a bit has crossed the doubleword boundary.
The new high half is \texttt{4000000000000000}. Inspect that transfer before
choosing a termination test for the complete exponent scan. The experiment also
reduces one modulo one to zero, making the zero-exponent identity explicit.

\subsection*{Check the modulus before following other pointers}

For every modular multiword entry, first read its four-limb modulus and determine
whether all limbs are zero. On zero, return status one before reading other caller
inputs or accessing output. Those other pointers can be deliberately invalid in
a test. Checking only the low modulus limb would incorrectly reject a valid
modulus such as $2^{32}$.

For positive modulus one, the result is zero, including an exponent-zero power.
For a positive larger modulus and exponent zero, the result is one. These are
successful arithmetic results, distinct from the zero-modulus error.

The scalar operands of \texttt{rsa\_pow32} are values, not addresses; its output
pointer still has the no-access promise on zero modulus. The gcd entry has no
modulus error path. Write each function's arguments and return convention beside
your test cases rather than applying one function's rules to all seven.

\begin{figure}[!ht]
  \centering
  \begin{tikzpicture}[alt={Actual native Memory row after zero-modulus rejection. All sixteen selected output bytes remain a5 and are outlined. The separately checked status register is outside this crop.}]
    \node[anchor=north west,inner sep=0,font=\small\bfseries] (title) at (0,0) {Zero modulus leaves the output unchanged};
    \node[anchor=north west,inner sep=0] (a) at ([yshift=-2mm]title.south west) {\MARSRevisionCrop{revision/captures/rsa/rsa-invalid.png}{0.99\linewidth}{0.03125}{0.33888889}{0.58203125}{0.60888889}};
    \begin{scope}[shift={(a.south west)},x={($(a.south east)-(a.south west)$)},y={($(a.north west)-(a.south west)$)}]
      \draw[black,line width=.7pt] (.20,.01) rectangle (.965,.467);
    \end{scope}
  \end{tikzpicture}
  \caption{All sixteen selected modular-output bytes remain a5 after rejecting modulus zero. The outline encloses the selected interval. Separate checks verify status one and no caller writes; the status register is not pictured.}
  \label{fig:revision-rsa-native-invalid}
\end{figure}


Success writes exactly four bytes for pow32, sixteen for modular multiword
results and 32 for the full product. Failure preserves all caller objects.
An unchanged final value alone does not prove that an illegal write never occurred:
a store of the same byte is still a write. The public checks observe access
intervals as well as final bytes.

In \texttt{worked/validation.asm}, the modulus is readable but zero, and
the other pointer arguments contain the invalid address one. The helper reads
all four modulus limbs and returns status one. Main then finds the separate
output word unchanged. The independent checks observe caller accesses as well
as bytes: an unchanged result alone would not prove that no forbidden access
occurred. Keep caller publication outside the completed callee's ownership
boundary.

\subsection*{Use a small RSA example as an arithmetic test}

The last native assessment groups exercise the actual caller and different
operations within one simulator. The repeated-power case feeds one result into
the inverse calculation and then repeats the original request. State left over
from a previous call must not replace current arguments. After completing all
entries, run the supplied mode-zero caller to termination, then use the final
assessment to recheck every interface on the completed source. Later edits can
invalidate an earlier pass. Choose another input and explain its expected result
from the arithmetic invariant before testing it.

For optional standalone Debug observations, keep \texttt{mode} at its supplied
default zero in the source passed to the preparation tool. To observe the scalar power in a fresh source from the book
checkout, use:

\Needspace{6\baselineskip}
\begin{Verbatim}[fontsize=\small]
python3 packages/027-rsa-arithmetic/prepare.py s06-1 \
  --source revision/labs/files/rsa-arithmetic.asm \
  --output /private/tmp/rsa-word-observation.asm
\end{Verbatim}

\Needspace{9\baselineskip}
Replace the source path with your completed working source when testing your
implementation. The tool installs data and mode in the output copy, refuses an
existing output path and retains your instruction text. Recipes
\texttt{s11-1}, \texttt{s18-1}, \texttt{s29-1} and \texttt{s36-1} select the
full product, carry-sensitive addition, modular product and full-width power
examples. An unfinished entry still stops with \texttt{break}; prepared inputs
do not supply its implementation.

For a separate scalar reverse run, set the first words of \texttt{left} and
\texttt{right} to 2790 and 2753, keep modulus 3233 and select mode one.
The full-width power caller uses mode six and reads its four-limb exponent from
\texttt{exponent} instead of \texttt{right}. The function's argument registers,
not the fixture labels themselves, determine which values it receives.

For standalone CLI checking, use the nine original public scripts listed in
\texttt{packages/027-rsa-arithmetic/package.json}. Their names separate gcd,
word power, product, addition, reduction, modular product and full-width power,
plus caller and repeated-call behavior. Use
\texttt{java -jar /path/to/cli.jar grade CHECK SOURCE}. Keep the supplied caller
when using its checks; private function labels and exact instruction counts are
not part of free-writing acceptance.

Build an independent test set that includes zero and one moduli, zero exponents,
top-bit-set scalar products, carries and borrows across every limb, maximum
full products, high exponent bits, exact addmod aliases and moved input/output
slices. Compare values with an arbitrary-precision calculation, then check
preservation and ownership separately. A round trip alone can hide two matching
mistakes.

The example explains RSA's modular-power arithmetic. A complete encryption or
signature scheme additionally defines how messages are encoded and checked.
RFC 8017 describes OAEP for an encryption scheme and PSS for a signature scheme;
this exercise implements neither encoding. Its four-limb numbers are chosen so
we can inspect the arithmetic, not as a usable RSA key size. Nor do the data-dependent
loops establish constant-time execution. Keep those scheme and implementation
questions separate from the numerical and memory contracts you have tested here.

